\chapter{Denotational Normal Forms}\label{ch08:chap}

\section{The canonicalization problem}\label{ch08:sec:problem}

Chapters~\ref{ch04:main} and~\ref{ch05:main} equipped the kernel with typed
composition and with powers where powers are legitimate. This chapter solves
the remaining Gate~1 problem: \emph{canonical names}. Given a finite set of
declared generator arrows, many different composition words produce the same
arrow. If two runs of an audit pipeline describe the same denotation by two
different words --- or worse, if the \emph{same} run does so depending on
interpreter hash order --- then downstream digests, caches, and comparisons
silently diverge. What is needed is a function $\NF_G$ that assigns to every
reachable denotation one distinguished word, such that:

\begin{enumerate}
\item the assignment is \emph{deterministic}: it depends only on the declared
  catalog (objects, generators, and their declared orders), never on
  interpreter state;
\item the assignment is \emph{denotational}: two words receive the same
  normal form exactly when they denote the same typed arrow, with equality of
  arrows being the extensional equality of
  Chapter~\ref{ch04:main} (same source object, same target object, same
  mapping tuple);
\item the assignment is \emph{total or absent}: either the complete table of
  reachable denotations is computed, or a resource exception is raised; a
  partial table is never presented as canonical.
\end{enumerate}

The solution adopted by the kernel, implemented in
\texttt{src/tks\_v8\_core/normal\_forms.py} and specified in the governing
Gate~1 sources, is \emph{shortlex minimization over a completed finite
closure}: explore all well-typed generator words in shortlex order, record
the first word achieving each denotation, and prove that first-discovery
under this exploration order coincides with shortlex-leastness. The
mathematical content of the chapter is exactly that proof, together with the
order theory it needs, the finiteness and complexity analysis, the resource
semantics, and two permanent boundaries that fence off what the construction
does \emph{not} claim.

Evidential status, stated up front and repeated where relevant: the
correctness results of this chapter (shortlex well-ordering, pruning
soundness, closure completeness) carry the ledger status of obligation
PO-TKS8-CORE-06 --- \textsc{informal correctness argument / formalization
open}. The arguments below are those informal arguments, in fully expanded
form. No Lean formalization of shortlex leastness exists; the pinned formal
seed's verified scope does not include this chapter's results
\cite{lean4}. The executable suite exercises normal-form correctness and
resource regressions on bounded families in the small-scope
tradition~\cite{jackson}; passing them establishes behavior on the covered
cases only. Everything is Ring~4, draft, not licensed.

\section{Generator catalogs}\label{ch08:sec:catalogs}

\begin{definition}[Generator]\label{ch08:def:generator}
A \emph{generator} is a pair $g=(\ell_g, f_g)$ of a non-blank string label
$\ell_g$ and a typed arrow $f_g$ of the finite-function category. The
implementation type is the frozen dataclass \texttt{Generator}, which
validates both components at construction.
\end{definition}

\begin{definition}[Catalog]\label{ch08:def:catalog}
A \emph{generator catalog} is a triple
\[
\cC=(\,\mathcal{O},\ G,\ \text{declared orders}\,)
\]
where $\mathcal{O}=(O_0,\dots,O_{m-1})$ is a finite, explicitly ordered
sequence of objects and $G=(g_0,\dots,g_{k-1})$ is a finite, explicitly
ordered sequence of generators, subject to the validation rules of
Invariant~\ref{ch08:inv:catalog}. The position of a generator in the
declared sequence is its \emph{rank}: $\operatorname{rank}(g_j)=j$. The
position of an object is its \emph{object rank}. Both orders are part of the
catalog's identity.
\end{definition}

\begin{invariant}[Catalog validation, as implemented]\label{ch08:inv:catalog}
The constructor of \texttt{NormalFormCatalog} enforces, raising
\texttt{ValidationError} on violation:
\begin{enumerate}
\item the object sequence is nonempty, and both sequences pass the
  ordered-input discipline (sets, mappings, and one-shot iterables are
  rejected, for the reasons of Chapter 3: order is meaning);
\item every declared object is a \texttt{FiniteObject} and every declared
  generator is a \texttt{Generator};
\item the objects are pairwise distinct, \emph{and} their names are pairwise
  distinct (name-inclusive identity would otherwise allow two objects that
  collide in error messages and reports);
\item the generator labels are pairwise distinct (labels are the alphabet of
  words; a duplicated label would make rank ill-defined);
\item every generator's arrow has both its source and its target among the
  declared objects (no arrow may smuggle in an undeclared object);
\item the resource caps \texttt{max\_arrows} and \texttt{max\_frontier} are
  positive integers (Booleans are explicitly rejected), with
  \texttt{max\_frontier} defaulting to \texttt{max\_arrows} and
  \texttt{max\_arrows} defaulting to $100{,}000$.
\end{enumerate}
Note what is \emph{not} required: generator arrows need not be distinct.
Two generators may carry extensionally equal arrows under different labels;
the theory below handles this case correctly (the two length-one words then
denote the same arrow, and shortlex picks the lower-ranked label).
\end{invariant}

\section{Typed words and denotations}\label{ch08:sec:words}

\begin{definition}[Typed word]\label{ch08:def:word}
A \emph{typed word} over a catalog $\cC$ is a pair $w=(S_w, \lambda_w)$ of a
\emph{start object} $S_w\in\mathcal{O}$ and a finite (possibly empty)
sequence of generator labels $\lambda_w=(\ell_1,\dots,\ell_n)$. The word is
\emph{well-typed} when the arrows compose in sequence from the start:
setting $X_0=S_w$, we require, for each $i$ from $1$ to $n$, that the
generator $g_{\ell_i}$ named by $\ell_i$ exists in the catalog and that its
declared source equals $X_{i-1}$; then $X_i$ is its declared target. The
implementation type \texttt{TypedWord} validates \emph{syntax} at
construction (a genuine \texttt{FiniteObject} start, an ordered tuple of
non-blank string labels); well-typedness is relative to a catalog and is
certified by \texttt{NormalFormCatalog.evaluate}, which raises
\texttt{ValidationError} on an unknown label or an undeclared start object
and raises the composition type error of Chapter~\ref{ch04:main} on a
junction mismatch.
\end{definition}

\begin{definition}[Denotation]\label{ch08:def:denotation}
The \emph{denotation} of a well-typed word $w=(S_w,(\ell_1,\dots,\ell_n))$
is the typed arrow
\[
\llbracket w\rrbracket \;=\; f_{\ell_n}\circ\cdots\circ f_{\ell_1}\circ
\id_{S_w}\;:\;S_w\longrightarrow X_n ,
\]
i.e.\ the composite of the generator arrows applied \emph{in reading
order}: the first label acts first. The empty word at $S$ denotes
$\id_S$. Equality of denotations is extensional equality of arrows: equal
source object, equal target object, equal mapping tuple. The implementation
computes $\llbracket w\rrbracket$ by a left fold, starting from
\texttt{identity(word.start)} and composing each successive generator arrow
on the left; the result is repackaged under a presentation name
(\texttt{word-evaluation[$n$]}) which, per the extensional-equality
discipline, never affects arrow equality. The \emph{denotation key} of an
arrow --- the triple (source object, target object, mapping tuple) --- is
the dictionary key under which the closure stores it.
\end{definition}

Two small but load-bearing observations about
Definition~\ref{ch08:def:denotation}:

\begin{remark}[Words determine their type path]\label{ch08:rem:typepath}
A nonempty well-typed word's start object is redundant data: it must equal
the declared source of its first generator. The start object earns its place
in the data type for the sake of the \emph{empty} words, one per declared
object, which are the only words whose denotation cannot be recovered from
labels alone. This is why identities are always in the closure --- they are
seeded as empty words --- even when no product of generators equals an
identity. It is also why, as we will use repeatedly below, \emph{two words
with the same denotation have the same start object}: the denotation's
source is the word's start.
\end{remark}

\begin{remark}[Reading order versus application order]\label{ch08:rem:order}
The word $(\ell_1,\ell_2)$ denotes $f_{\ell_2}\circ f_{\ell_1}$: reading
order is application order, which is the \emph{reverse} of the conventional
right-to-left composition notation. The convention is chosen so that
extending a word on the right corresponds to composing on the left with a
new generator --- exactly the step the closure algorithm performs --- and so
that a word reads like an execution trace. Nothing depends on the
convention except consistency, which the implementation enforces in one
place (\texttt{evaluate}).
\end{remark}

\section{The shortlex order}\label{ch08:sec:shortlex}

\begin{definition}[Shortlex]\label{ch08:def:shortlex}
Fix a catalog $\cC$. For a typed word $w$, let the \emph{rank sequence} of
$w$ be the sequence of generator ranks of its labels, and let its
\emph{key} be the triple
\[
\kappa(w)=\bigl(\,|w|,\ \text{rank sequence of } w,\
\text{object rank of } S_w\,\bigr),
\]
where $|w|$ is the number of labels. Keys are compared lexicographically:
first by length; then by rank sequence in lexicographic order (compare
position by position; the first differing position decides); then by start
object rank. Write $u<_{\mathrm{slex}}v$ iff $\kappa(u)<\kappa(v)$, and
$u\le_{\mathrm{slex}}v$ iff $\kappa(u)\le\kappa(v)$. This is precisely the
key computed by the implementation method \texttt{shortlex\_key}.
\end{definition}

The third component deserves comment before the theory. For nonempty words
the first two components already determine the word: the rank sequence
determines the label sequence (ranks and labels are in bijection by
Invariant~\ref{ch08:inv:catalog}(4)), and the label sequence determines the
start object (Remark~\ref{ch08:rem:typepath}). Only the empty words at
distinct objects share the key prefix $(0,())$, and the object-rank
component separates them. Thus:

\begin{lemma}[Keys are injective]\label{ch08:lem:keyinjective}
Distinct typed words over the same catalog have distinct keys. Consequently
$<_{\mathrm{slex}}$ is a strict total order on typed words: for any two
distinct words, exactly one of $u<_{\mathrm{slex}}v$, $v<_{\mathrm{slex}}u$
holds.
\end{lemma}

\begin{proof}
Suppose $\kappa(u)=\kappa(v)$. Then $|u|=|v|$ and the rank sequences agree.
If the common length is positive, the rank sequences determine equal label
sequences, whose first labels name generators with equal declared sources,
so the start objects agree and $u=v$. If the common length is zero, both
words are empty and the equal third components force equal start objects
(object ranks are positions in a sequence of \emph{distinct} objects), so
again $u=v$. Totality and trichotomy now follow from totality of the
lexicographic comparison of integer triples-of-sequences: lengths are
integers; rank sequences of equal length are compared position by position
over the total order of integers; object ranks are integers.
\end{proof}

\begin{lemma}[Well-foundedness; shortlex is a well-order]\label{ch08:lem:wellfounded}
Every nonempty set $W$ of typed words over a catalog with $k$ generators
and $m$ objects has a unique $<_{\mathrm{slex}}$-least element. Equivalently,
there is no infinite strictly $<_{\mathrm{slex}}$-descending sequence of
words.
\end{lemma}

\begin{proof}
Let $L=\min\{|w| : w\in W\}$, which exists because lengths are natural
numbers. The set $W_L$ of words in $W$ of length $L$ is nonempty and
finite: there are at most $k^L$ label sequences of length $L$ and at most
$m$ start objects, so $|W_L|\le m\,k^L$ (for $L\ge 1$ even $\le k^L$, by
Remark~\ref{ch08:rem:typepath}). A finite nonempty set under a strict total
order has a unique least element; the least element of $W_L$ is
$<_{\mathrm{slex}}$-below every other member of $W$ (shorter beats longer;
within $W_L$ it is least by construction). For the equivalent formulation:
an infinite strictly descending sequence would have nonincreasing lengths,
hence eventually constant length $L'$, and would then descend strictly and
infinitely inside the finite set of words of length $L'$ from at most $m$
starts --- impossible.
\end{proof}

\begin{lemma}[Right-append monotonicity]\label{ch08:lem:append}
Let $u,u'$ be typed words with the same start object and
$u<_{\mathrm{slex}}u'$, and let $s$ be a label sequence. Write $us$ for the
word $(S_u,\ \lambda_u\!\cdot\! s)$ and $u's$ for
$(S_{u'},\ \lambda_{u'}\!\cdot\! s)$. Then $us<_{\mathrm{slex}}u's$.
\end{lemma}

\begin{proof}
Case 1: $|u|<|u'|$. Then $|us|=|u|+|s|<|u'|+|s|=|u's|$, and shorter wins.
Case 2: $|u|=|u'|$. Since the start objects agree, the key comparison was
decided by the rank sequences: $\operatorname{ranks}(u)$ precedes
$\operatorname{ranks}(u')$ lexicographically, so they differ at some first
position $i\le|u|$, with the $u$-entry smaller. The rank sequences of $us$
and $u's$ have equal length $|u|+|s|$ and agree with the originals on the
first $|u|$ positions; their first difference is therefore still at position
$i$, still in $u$'s favor. Hence $us<_{\mathrm{slex}}u's$. (The third key
component is equal on both sides and never consulted.)
\end{proof}

\begin{remark}[Why shortlex and not plain lex]\label{ch08:rem:whynotlex}
The length-first component is not a stylistic preference; it is what makes
the order usable at all. Consider plain lexicographic order on words, with
the standard convention that a proper prefix precedes its extensions. Under
that order, minimization over an infinite language can fail to make sense
algorithmically even when minima exist, and --- decisive for us ---
exploration order and well-foundedness part company. Concretely, with two
generators $a<b$, the set $\{b,\ ab,\ aab,\ aaab,\ \dots\}$ is an infinite
strictly \emph{descending} chain in plain lex ($aab<ab<b$, and each added
$a$ descends further), so plain lex is not well-founded on words and
``the lex-least word denoting $a$'' need not exist: a denotation reachable
by all of $b, ab, aab, \dots$ (easily arranged when $a$'s arrow acts as an
identity on the relevant object) has no least representative at all.
Shortlex repairs this with the length component, which caps every
descending chain (Lemma~\ref{ch08:lem:wellfounded}) and simultaneously
aligns the order with breadth-first exploration: all words of length $n$
precede all words of length $n+1$, so a frontier processed in shortlex
order is a breadth-first frontier with a deterministic within-level order.
The alignment is what lets one algorithm serve as both an exhaustive
explorer and a minimizer --- the content of
Theorem~\ref{ch08:thm:soundness}.
\end{remark}

\section{Normal forms}\label{ch08:sec:nf}

\begin{definition}[Reachable denotation; $\NF_G$]\label{ch08:def:nf}
An arrow $a$ is a \emph{reachable denotation} of the catalog $\cC$ when
$a=\llbracket w\rrbracket$ for some well-typed word $w$ over $\cC$ (the
empty words make every $\id_{O_i}$ reachable). For reachable $a$, the
\emph{normal form} of $a$ is
\[
\NF_G(a)\;=\;\min_{<_{\mathrm{slex}}}\ \bigl\{\, w \text{ well-typed over }
\cC \;\bigm|\; \llbracket w\rrbracket = a \,\bigr\}.
\]
\end{definition}

\begin{proposition}[$\NF_G$ is well defined]\label{ch08:prop:nfwelldef}
For every reachable denotation $a$, the set of well-typed words denoting
$a$ is nonempty, and it has a unique shortlex-least element.
\end{proposition}

\begin{proof}
Nonemptiness is the definition of reachability; existence and uniqueness of
the least element are Lemma~\ref{ch08:lem:wellfounded} (uniqueness also
directly from Lemma~\ref{ch08:lem:keyinjective}: two distinct least elements
would be shortlex-comparable, contradicting joint minimality).
\end{proof}

The normal form is the object the rest of the kernel consumes: the method
\texttt{normalize(word)} evaluates a word and returns
$\NF_G(\llbracket w\rrbracket)$; the method
\texttt{representative\_for(arrow)} looks an arrow up by denotation key,
raising \texttt{UnreachableDenotationError} when the arrow is not in the
generated closure; the property \texttt{normal\_forms} returns all
representatives sorted by shortlex key; and \texttt{denotation\_count}
reports the closure size. What remains is to \emph{compute} $\NF_G$ --- and
to prove that the computation is right.

\section{The representative-first closure algorithm}\label{ch08:sec:algorithm}

\begin{algorithm}[Representative-first closure]\label{ch08:alg:closure}
Input: a validated catalog $\cC$ with caps \texttt{max\_arrows},
\texttt{max\_frontier}. State: a min-heap (priority queue) \emph{frontier}
of entries $(\kappa(w),\ \text{serial},\ w,\ \llbracket w\rrbracket)$
ordered by key $\kappa$ with a monotone serial number as tiebreaker; a set
\emph{queued} of denotation keys currently represented in the frontier; a
table \emph{reps} from denotation keys to words (with a parallel table of
the arrows themselves). Steps:
\begin{enumerate}
\item \textbf{Seed.} For each declared object $O$ in declared order: if
  $|\text{queued}|\ge{}$\texttt{max\_frontier}, raise
  \texttt{ClosureLimitExceeded}; otherwise push the empty word at $O$ with
  its denotation $\id_O$ and add $\id_O$'s key to \emph{queued}.
\item \textbf{Loop.} While the frontier is nonempty:
  \begin{enumerate}
  \item Pop the entry with least key; let it be $(w,a)$ with denotation key
    $d$. Remove $d$ from \emph{queued}.
  \item \textbf{Duplicate check.} If $d\in\text{reps}$, discard $(w,a)$ and
    continue the loop (the word is recorded as a duplicate and \emph{not
    expanded}).
  \item \textbf{Cap check.} If $|\text{reps}|\ge{}$\texttt{max\_arrows},
    raise \texttt{ClosureLimitExceeded}.
  \item \textbf{Accept.} Record $\text{reps}[d]=w$.
  \item \textbf{Expand.} For each generator $g$ in declared (rank) order
    whose declared source equals the target of $a$: form the extension
    $w'=w\cdot\ell_g$ and its denotation $a'=f_g\circ a$ with key $d'$. If
    $d'\in\text{reps}$ or $d'\in\text{queued}$, skip (denotation
    deduplication). Otherwise: if $|\text{queued}|\ge{}$%
    \texttt{max\_frontier}, raise \texttt{ClosureLimitExceeded}; else push
    $(w',a')$ and add $d'$ to \emph{queued}.
  \end{enumerate}
\item \textbf{Return.} When the frontier empties without a cap being hit,
  the table \emph{reps} is complete and the catalog object becomes usable.
\end{enumerate}
This is a faithful transcription of \texttt{NormalFormCatalog.\_build}. The
serial number exists so that heap entries are always comparable even if two
keys were equal; by Lemma~\ref{ch08:lem:keyinjective} keys of distinct words
are never equal, so the serial is a belt-and-braces guard for the runtime,
not a mathematical ingredient.
\end{algorithm}

Three design points before the proofs. First, \emph{representative-first}
names the policy of step 2(b): only the first word to arrive at a denotation
is ever expanded; all later arrivals are pruned unexpanded. The entire
correctness question is whether this pruning can ever lose a smaller
representative or an entire denotation --- Section~\ref{ch08:sec:soundness}
proves it cannot. Second, the \emph{queued} set prunes even earlier: a
candidate whose denotation is already awaiting processing in the frontier is
not pushed at all. This is an optimization with teeth --- it bounds the
frontier by the number of distinct pending denotations --- and it carries a
proof obligation of its own (Lemma~\ref{ch08:lem:skip}), because dropping a
candidate is only sound if the word already queued for that denotation is at
least as good. Third, the caps are checked \emph{before} accepting or
pushing, so the exception fires instead of the table silently saturating;
the construction never returns a partial table
(Section~\ref{ch08:sec:resources}).

\subsection{The closure as a graph search}\label{ch08:sec:graphview}

A reader from combinatorial semigroup theory will recognize the skeleton of
Algorithm~\ref{ch08:alg:closure}. Define the \emph{typed denotation graph}
of a catalog: vertices are the reachable denotations; there is a labeled
edge $a\xrightarrow{\ \ell\ }a'$ whenever the generator labeled $\ell$ has
source equal to the target of $a$ and $a'=f_\ell\circ a$. This is the
typed, multi-object analogue of the (left) Cayley graph of a monoid with
respect to a generating set \cite{pin,howie}, with the identities $\id_O$
playing the role of the neutral vertex --- one per declared object, since
the ``monoid'' here is really the arrow family of a category. Under this
reading:
\begin{itemize}
\item the closure is a graph search from the $m$ seed vertices;
\item a representative $\NF_G(a)$ is a \emph{geodesic} label sequence from
  a seed to $a$ --- shortest first, by the length component --- with ties
  broken deterministically by the rank-lexicographic component;
\item the duplicate check of step 2(b) is the classic ``visited'' set of
  breadth-first search; the \emph{queued} set of step 2(e) is the equally
  classic ``discovered but not yet expanded'' refinement that prevents the
  frontier from holding one vertex twice;
\item Theorem~\ref{ch08:thm:soundness} is then the typed, tie-broken
  sharpening of the standard fact that breadth-first search computes
  shortest paths: not merely \emph{some} geodesic per vertex, but the
  lexicographically least one, which is what determinism requires.
\end{itemize}
The analogy also delimits itself. Cayley-graph methods in semigroup theory
typically continue toward presentations, relations, and rewriting
\cite{howie}; the kernel stops deliberately short of all three
(Boundary~\ref{ch08:bd:norewrite}). The graph view is offered as an aid to
intuition and a bridge to the literature, not as an alternative
formalization: the definitions that govern are those of
Sections~\ref{ch08:sec:catalogs}--\ref{ch08:sec:algorithm}.

\section{Soundness of pruning}\label{ch08:sec:soundness}

Throughout this section, fix a validated catalog and consider a run of
Algorithm~\ref{ch08:alg:closure} that terminates normally (no cap raised).
``Popped'' means returned by step 2(a); the popped words form a finite
sequence $p_1,p_2,\dots,p_T$ in temporal order.

The proof rests on the two invariants announced in the governing sources:
the \emph{dequeue-order} invariant (words leave the frontier in shortlex
order) and the \emph{prefix-domination} invariant (a better prefix stays
better under any common suffix). We prove each, plus the auxiliary
skip-domination lemma that the frontier deduplication requires.

\begin{lemma}[Invariant I: dequeue monotonicity]\label{ch08:lem:dequeue}
$\kappa(p_1)\le\kappa(p_2)\le\cdots\le\kappa(p_T)$; that is, words are
popped in nondecreasing (by Lemma~\ref{ch08:lem:keyinjective}, strictly
increasing) shortlex order.
\end{lemma}

\begin{proof}
Every non-seed frontier entry is pushed during the expansion step of some
pop, and an extension $p_t\cdot\ell$ satisfies
$\kappa(p_t\cdot\ell)>\kappa(p_t)$ because its length is $|p_t|+1$. Now fix
$t$ and compare $p_{t+1}$ with $p_t$. At the moment $p_t$ was popped, the
entry that will later be popped as $p_{t+1}$ either (i) was already in the
frontier --- in which case $\kappa(p_{t+1})\ge\kappa(p_t)$ because the heap
returned $p_t$ as the \emph{minimum} of the then-current frontier --- or
(ii) had not yet been pushed, in which case it is pushed during the
expansion of $p_t$ itself (pushes happen only inside expansions, and the
only expansion between the two pops is that of $p_t$), whence
$\kappa(p_{t+1})=\kappa(p_t\cdot\ell)>\kappa(p_t)$ for some label $\ell$.
Either way $\kappa(p_t)\le\kappa(p_{t+1})$, and induction over $t$ finishes
the claim. (Seeds are all pushed before the first pop, so case (ii) never
involves a seed.)
\end{proof}

\begin{lemma}[Invariant II: prefix domination under common suffixes]\label{ch08:lem:prefix}
Let $u,u'$ be well-typed words with
$\llbracket u\rrbracket=\llbracket u'\rrbracket$ and
$u<_{\mathrm{slex}}u'$. Then for every label sequence $s$ such that $u's$ is
well-typed: $us$ is well-typed,
$\llbracket us\rrbracket=\llbracket u's\rrbracket$, and
$us<_{\mathrm{slex}}u's$.
\end{lemma}

\begin{proof}
Equal denotations have equal sources and equal targets, so $u$ and $u'$
share a start object (Remark~\ref{ch08:rem:typepath}) and end at the same
object $Z$. Well-typedness of an appended suffix depends only on the object
at the junction and on the suffix's own internal typing: $u's$ well-typed
means the first label of $s$ names a generator with source $Z$ and the
remaining labels chain on from there; the identical condition makes $us$
well-typed. For the denotation, evaluation is a fold, so
$\llbracket us\rrbracket = F_s\circ\llbracket u\rrbracket$ where $F_s$
denotes the composite of the suffix's generator arrows (an arrow
$Z\to Z'$); likewise
$\llbracket u's\rrbracket=F_s\circ\llbracket u'\rrbracket$; the right-hand
factors are equal by hypothesis, so the composites are extensionally equal.
The order claim is Lemma~\ref{ch08:lem:append}, whose same-start hypothesis
we just established.
\end{proof}

\begin{lemma}[Skip domination]\label{ch08:lem:skip}
Suppose that during the expansion of a popped word $p_t$, the candidate
$w'=p_t\cdot\ell$ with denotation key $d'$ is \emph{skipped} in step 2(e).
Then at that moment there exists a well-typed word $q$ with denotation key
$d'$ and $q<_{\mathrm{slex}}w'$, where $q$ is either the accepted
representative of $d'$ (if $d'\in\text{reps}$) or the unique frontier word
with denotation key $d'$ (if $d'\in\text{queued}$).
\end{lemma}

\begin{proof}
\emph{Accepted case.} The representative $q=\text{reps}[d']$ was popped as
some $p_s$ with $s\le t$, so $\kappa(q)\le\kappa(p_t)<\kappa(w')$ by
Lemma~\ref{ch08:lem:dequeue} and the strict length increase of extensions.

\emph{Queued case.} The frontier word $q$ with key $d'$ was pushed either as
a seed or during the expansion of some earlier pop $p_s$ with $s\le t$.
Since $q$ and $w'$ have the same denotation, they have the same start object
(Remark~\ref{ch08:rem:typepath}); note this rules out the one scenario in
which the object-rank tiebreaker could interfere, namely distinct empty-word
seeds --- two seeds have different denotations (identities of distinct
objects), so if $q$ is a seed then $w'$, being nonempty, is longer and
$q<_{\mathrm{slex}}w'$ immediately. Otherwise $q=p_s\cdot\ell'$ for some
$s\le t$ and label $\ell'$. Compare:
\begin{itemize}
\item If $s<t$ and $|p_s|<|p_t|$: then $|q|<|w'|$ and $q<_{\mathrm{slex}}w'$.
\item If $s<t$ and $|p_s|=|p_t|$: by Lemma~\ref{ch08:lem:dequeue} and key
  injectivity, $\kappa(p_s)<\kappa(p_t)$ with equal lengths, so the rank
  sequences differ at some first position inside the common length, in
  $p_s$'s favor; appending one label to each does not disturb that first
  difference, so $q<_{\mathrm{slex}}w'$ (the argument of
  Lemma~\ref{ch08:lem:append}, applied across two different prefixes).
\item If $s=t$: both $q$ and $w'$ extend the same popped word $p_t$, and
  generators are enumerated in declared rank order, so the earlier push has
  the smaller final rank: $\operatorname{rank}(\ell')<\operatorname{rank}(\ell)$,
  whence $q<_{\mathrm{slex}}w'$ (equal lengths, equal prefix, smaller last
  entry).
\end{itemize}
In every case $q<_{\mathrm{slex}}w'$, as claimed.
\end{proof}

With the three lemmas in hand, we can state and prove the main result. Call
a word $w$ \emph{hereditarily minimal} when every prefix of $w$ (including
$w$ itself, with its start object) is the shortlex-least word for that
prefix's own denotation.

\begin{lemma}[Prefixes of normal forms are normal forms]\label{ch08:lem:hereditary}
If $w=\NF_G(a)$ for some reachable $a$, then $w$ is hereditarily minimal.
\end{lemma}

\begin{proof}
Let $p$ be a proper prefix of $w$, say $w=ps$ with $s$ nonempty, and
suppose toward a contradiction that some well-typed $r$ has
$\llbracket r\rrbracket=\llbracket p\rrbracket$ and $r<_{\mathrm{slex}}p$.
By Lemma~\ref{ch08:lem:prefix} (with $u=r$, $u'=p$, suffix $s$), the word
$rs$ is well-typed, denotes $\llbracket ps\rrbracket=a$, and satisfies
$rs<_{\mathrm{slex}}ps=w$ --- contradicting $w=\NF_G(a)$. The case $p=w$ is
the definition of $\NF_G$.
\end{proof}

\begin{theorem}[Closure correctness; \textnormal{status: informal correctness argument, formalization open (PO-TKS8-CORE-06)}]\label{ch08:thm:soundness}
Suppose Algorithm~\ref{ch08:alg:closure} terminates normally on a validated
catalog. Then the computed table \emph{reps} satisfies:
\begin{enumerate}
\item (\emph{Completeness}) the key set of \emph{reps} is exactly the set of
  reachable denotations of the catalog;
\item (\emph{Minimality}) for every reachable $a$,
  $\text{reps}[a]=\NF_G(a)$.
\end{enumerate}
\end{theorem}

\begin{proof}
We first prove the following claim by strong induction on the shortlex
order (legitimate by Lemma~\ref{ch08:lem:wellfounded}):

\medskip
\noindent\emph{Claim.} For every reachable denotation $a$, the word
$w^\ast=\NF_G(a)$ is pushed at some point of the run, and when it is popped,
$a$ is not yet in \emph{reps}; consequently $a$ enters \emph{reps} with
representative $w^\ast$.

\medskip
\noindent\emph{Proof of claim.} Fix $a$ and $w^\ast=\NF_G(a)$, and assume
the claim holds for every reachable denotation whose normal form is
shortlex-smaller than $w^\ast$.

\emph{Getting $w^\ast$ pushed.} If $w^\ast$ is empty, it is a seed and is
pushed in step 1. Otherwise write $w^\ast=p\cdot\ell$ where $p$ is the
longest proper prefix. By Lemma~\ref{ch08:lem:hereditary},
$p=\NF_G(\llbracket p\rrbracket)$, and $p<_{\mathrm{slex}}w^\ast$, so the
induction hypothesis applies to $\llbracket p\rrbracket$: the word $p$ is
popped and accepted as the representative of its denotation. During the
expansion of $p$, the candidate $p\cdot\ell=w^\ast$ is generated (the
generator named $\ell$ has source equal to the target of
$\llbracket p\rrbracket$ because $w^\ast$ is well-typed). Could it be
skipped? By Lemma~\ref{ch08:lem:skip}, a skip would produce a word
$q<_{\mathrm{slex}}w^\ast$ with $\llbracket q\rrbracket=a$ ---
contradicting minimality of $w^\ast$. So $w^\ast$ is pushed.

\emph{Popping $w^\ast$ first among $a$'s words.} The run terminates
normally with an empty frontier, and every popped entry is either accepted
or discarded as a duplicate, so every pushed entry is eventually popped; in
particular $w^\ast$ is popped, say as $p_t$. Suppose $a\in\text{reps}$
already at that moment, via an earlier accepted word $r=p_s$, $s<t$, with
$\llbracket r\rrbracket=a$. By Lemma~\ref{ch08:lem:dequeue} and key
injectivity, $r\ne w^\ast$ implies $\kappa(r)<\kappa(w^\ast)$ --- again
contradicting $w^\ast=\NF_G(a)$. Hence $a\notin\text{reps}$ when $w^\ast$
is popped, the duplicate check passes, and (normal termination having
excluded the cap branch) $w^\ast$ is accepted as $\text{reps}[a]$. This
proves the claim.

\medskip
\emph{Completeness.} Every key in \emph{reps} is the denotation key of some
evaluated word, hence reachable. Conversely, by the claim every reachable
denotation is accepted. \emph{Minimality.} A denotation enters \emph{reps}
exactly once (the duplicate check forbids overwriting), and by the claim the
entering word is $\NF_G(a)$.
\end{proof}

\begin{corollary}[Determinism and stability]\label{ch08:cor:stability}
On normal termination, the table \emph{reps} --- keys, representatives, and
the sorted listing \texttt{normal\_forms} --- is a pure function of the
catalog identity: the object sequence (names, carriers, and order), the
generator sequence (labels, arrows, and order), and nothing else. Two runs
on equal catalogs produce equal tables, on any machine, under any
\texttt{PYTHONHASHSEED}.
\end{corollary}

\begin{proof}
Algorithm~\ref{ch08:alg:closure} consults only: declared object order
(seeding), declared generator order (expansion and ranks), arrow composition
(deterministic by Chapter~\ref{ch04:main}), and shortlex keys (functions of
the declared orders). The heap's behavior is determined by keys and the
deterministic serial tiebreaker; no iteration over a hash-ordered container
occurs on any code path (the ordered-input discipline of
Invariant~\ref{ch08:inv:catalog} exists to guarantee precisely this).
Theorem~\ref{ch08:thm:soundness} then pins the \emph{output} to the
mathematical object $\NF_G$, which depends only on the catalog.
\end{proof}

\section{Finiteness and complexity}\label{ch08:sec:complexity}

\begin{proposition}[Finiteness of the closure]\label{ch08:prop:finite}
Let the catalog declare objects $O_0,\dots,O_{m-1}$ with carrier sizes
$n_0,\dots,n_{m-1}$. The number of reachable denotations $N$ satisfies
\[
N \;\le\; \sum_{i=0}^{m-1}\sum_{j=0}^{m-1} n_j^{\,n_i},
\]
since for each ordered pair of objects $(O_i,O_j)$ there are at most
$n_j^{\,n_i}$ total maps $O_i\to O_j$, and every denotation is a total map
between declared objects. (The convention $0^0=1$ counts the unique empty
map, so empty carriers are covered.) Consequently
Algorithm~\ref{ch08:alg:closure} terminates: each denotation is accepted at
most once, each acceptance pushes at most $k$ extensions, at most one
frontier entry exists per pending denotation, and the frontier therefore
empties after finitely many pops.
\end{proposition}

\begin{proof}
The bound is the counting statement just made. For termination, observe that
every pop either accepts a new denotation or discards a duplicate; there are
at most $N$ acceptances; total pushes are bounded by $m$ seeds plus $k$ per
acceptance, i.e.\ $m+kN$; every pop removes one entry; so the loop performs
at most $m+kN$ pops and halts with an empty frontier (or raises earlier).
\end{proof}

\begin{remark}[Cost accounting]\label{ch08:rem:cost}
Write $N$ for the actual closure size, $k$ for the generator count, $m$ for
the object count, $n=\max_i n_i$, and $F$ for the peak frontier size. The
run performs at most $m+kN$ pushes and pops, each costing
$O(\log F)$ comparisons of keys; a key comparison costs $O(L)$ where $L$ is
the length of the longer word, and $L$ is bounded by $N$ itself in the worst
case (each accepted representative extends an earlier one by one label, and
representatives are pairwise distinct denotations, so a representative's
length never exceeds the number of acceptances that preceded it --- a chain
of new denotations). Each expansion performs at most $k$ compositions at
$O(n)$ integer work apiece plus one denotation-key construction at $O(n)$.
Total: $O\bigl((m+kN)\,(\,k\,n + L\log F\,)\bigr)$ in the worst case, with
$L\log F$ typically far below its bound because representative lengths in
practice grow like the diameter of the denotation graph, not like $N$.
\end{remark}

\begin{remark}[Frontier growth, and what deduplication buys]\label{ch08:rem:frontier}
Without the \emph{queued} set, the frontier could hold many words per
denotation --- up to one per accepted predecessor --- and its size could
approach $kN$. With deduplication, the frontier holds \emph{at most one
entry per pending denotation}, so $F\le N$ always, and
$F\le{}$\texttt{max\_frontier} is enforced besides. The price is the proof
obligation discharged by Lemma~\ref{ch08:lem:skip}: dropping a later word
for an already-queued denotation is only sound because the queued word is
strictly shortlex-smaller. The pattern --- an optimization that is also a
theorem --- recurs throughout the kernel, and the monograph's policy is to
state the theorem next to the optimization every time.
\end{remark}

\begin{boundary}[The bound is doubly exponential in carrier size, and that is fine]\label{ch08:bd:bound}
The worst-case closure bound $\sum_{i,j}n_j^{n_i}$ grows brutally: two
objects of size $6$ already admit $6^6=46{,}656$ maps per direction. This is
not a defect to engineer away; it is the true size of the function space,
and any ``canonical form'' scheme that returned answers faster than the
space's size on adversarial catalogs would be sampling, not closing. The
kernel's answer to the growth is not cleverness but honesty: dual resource
caps and a fail-closed exception, next.
\end{boundary}

\section{Resource semantics}\label{ch08:sec:resources}

\begin{construction}[Dual caps and fail-closed completion]\label{ch08:con:resources}
The constructor accepts \texttt{max\_arrows} (default $100{,}000$; bounds
the number of \emph{accepted} denotations) and \texttt{max\_frontier}
(default: equal to \texttt{max\_arrows}; bounds the number of \emph{pending}
denotations in the frontier at any moment). Three guard sites exist, all
raising \texttt{ClosureLimitExceeded}: seeding a frontier already at the
frontier cap; accepting a denotation when the table already holds
\texttt{max\_arrows} entries; and pushing an extension when the frontier is
at the cap. The exception's message states, in both cap cases, that
\emph{no partial normal-form catalog was returned} --- and none is: the
exception escapes the constructor, so no \texttt{NormalFormCatalog} object
exists afterward; there is no partially built catalog to misuse.
\end{construction}

\begin{ruleenv}[Fail-closed canonicalization]\label{ch08:rule:failclosed}
A normal-form table is either \emph{complete} --- every reachable
denotation present, every representative shortlex-least, per
Theorem~\ref{ch08:thm:soundness} --- or it is \emph{absent}. There is no
third state. In particular the implementation never: returns the accepted
portion of an interrupted closure; marks a table ``best effort''; or
degrades to a sampled or truncated table under resource pressure.
\end{ruleenv}

The rationale deserves spelling out, because the alternative is common
engineering practice. A truncated ``canonical'' table is worse than none: a
lookup that fails could mean \emph{unreachable} or could mean
\emph{unexplored}, and a representative drawn from a truncated exploration
may be beaten by an unexplored shorter word, so equality of representatives
would no longer certify anything. Every downstream consumer of the catalog
--- digest computation, cross-run comparison, deduplication --- relies on
representatives being genuine minima. The kernel therefore converts resource
exhaustion into a loud, typed failure at construction time, and the two caps
let a caller choose the failure threshold rather than whether failure is
loud. Separately, \texttt{representative\_for} distinguishes the
\emph{unreachable} case with its own exception,
\texttt{UnreachableDenotationError}: on a \emph{completed} table,
``not found'' soundly means ``not generated by the declared generators,''
and the exception name says so.

\begin{example}[Tripping the cap, concretely]\label{ch08:ex:captrip}
To see the guard fire on an honest workload rather than a contrived one,
declare a single object $T=(\texttt{Three},(0,1,2))$ and three generators on
it: $c$ with the $3$-cycle $0\mapsto 1\mapsto 2\mapsto 0$; $t$ with the
transposition swapping $0$ and $1$; and $e$ with the collapse
$0\mapsto 0$, $1\mapsto 0$, $2\mapsto 2$. It is a standard fact of
transformation-semigroup theory that a generating set of the symmetric
group together with one idempotent of rank $n-1$ generates the full
transformation monoid on $n$ letters \cite{howie}; for $n=3$ the reachable
denotations are therefore all $3^3=27$ endomaps of $T$ (plus nothing else:
one object, so every denotation is an endomap, and the empty word supplies
the identity). A catalog constructed over these generators with
\texttt{max\_arrows}${}=27$ or more completes with
\texttt{denotation\_count}${}=27$; the same catalog with
\texttt{max\_arrows}${}=10$ raises \texttt{ClosureLimitExceeded} during the
eleventh acceptance, and no catalog object exists afterward. The example
makes the design point tangible: the cap is not an error in the catalog ---
the closure is finite and perfectly computable --- but a caller-chosen
budget, and the failure it produces is loud, typed, and total, never a
$10$-entry table masquerading as the canonical closure of a $27$-element
monoid.
\end{example}

\section{Generator-order relativity}\label{ch08:sec:orderrel}

Canonical forms are relative to the declared generator order. This is not a
bug to be normalized away but a definitional fact to be documented and
enforced; this section proves the invariance that does hold, exhibits the
variance that does not, and states the resulting comparison discipline.

\begin{proposition}[Denotations are order-invariant; representatives are not]\label{ch08:prop:orderrel}
Let $\cC$ and $\cC'$ be catalogs with the same object set and the same
generator \emph{set} (same labels, same arrows) but different declared
generator orders. Then:
\begin{enumerate}
\item the reachable denotation sets of $\cC$ and $\cC'$ are equal;
\item the closure tables have equal key sets, and
  \texttt{denotation\_count} agrees;
\item the representative $\NF_G$ of a given denotation may differ between
  $\cC$ and $\cC'$; when it differs, the two representatives denote the
  same arrow.
\end{enumerate}
\end{proposition}

\begin{proof}
(1) Reachability quantifies over well-typed words, and well-typedness of a
word (Definition~\ref{ch08:def:word}) mentions labels, arrows, sources, and
targets --- never ranks. The same words are well-typed over both catalogs,
with the same denotations. (2) Immediate from (1) and completeness
(Theorem~\ref{ch08:thm:soundness}). (3) The normal form minimizes over the
same set of words but under a different order, so the minimum may move;
whatever word is selected still denotes the arrow being represented, by
definition of the set minimized over.
\end{proof}

\begin{example}[Worked order-relativity: two commuting transpositions]\label{ch08:ex:orderrel}
Declare one object $Q=(\texttt{Quad},(0,1,2,3))$ and two generators on it:
$a$ with arrow the transposition swapping positions $0$ and $1$ (fixing
$2,3$), and $b$ with arrow the transposition swapping positions $2$ and $3$
(fixing $0,1$). The two arrows commute: $a\circ b=b\circ a$ is the double
transposition $\tau$ swapping both pairs. The reachable denotations are the
four arrows of the Klein four-group: $\id_Q$, $a$, $b$, $\tau$ (products of
longer words fold back: $aa=bb=\id$, $ab=ba=\tau$, and so on).

Under declared order $(a,b)$ --- ranks $a\mapsto 0$, $b\mapsto 1$ --- the
closure accepts, in pop order: $()$ for $\id_Q$; $(a)$ for $a$; $(b)$ for
$b$; then among length-two words, $(a,a)$ duplicates $\id$, and
$(a,b)$ is popped before $(b,a)$ because rank sequence $(0,1)<(1,0)$
lexicographically --- so $\tau$ gets representative $(a,b)$, and when
$(b,a)$ is generated it is skipped (denotation already present).

Under declared order $(b,a)$ --- ranks $b\mapsto 0$, $a\mapsto 1$ --- the
identical word set is explored under the mirrored order, and $\tau$ gets
representative $(b,a)$: rank sequence $(0,1)$ again, but $(0,1)$ now
\emph{spells} ``$b$ then $a$.'' The denotation table is identical in both
runs --- four keys, the same four arrows --- while the representative of
$\tau$ differs as a word. Both representatives evaluate to $\tau$; neither
is wrong; each is canonical \emph{relative to its catalog}.
\end{example}

\begin{boundary}[Catalog identity is required for cross-comparison]\label{ch08:bd:catalogidentity}
A normal form is meaningful only as a pair (catalog identity, word).
Comparing representatives produced under different catalogs --- different
generator orders, different labels, different object declarations, or
different generator sets --- is meaningless: equal words may denote
different arrows (if labels were rebound) and equal arrows may receive
different words (Example~\ref{ch08:ex:orderrel}). Any workflow that stores,
transmits, or diffs normal forms must bind them to a versioned catalog
identity --- in the kernel's artifact discipline, the declared sequences and
their hashes --- and must treat a catalog change as a schema change that
invalidates all stored representatives, even when the denotation sets
coincide (Proposition~\ref{ch08:prop:orderrel} guarantees they often do,
which makes the silent-mismatch failure mode \emph{more} dangerous, not
less: the keys line up while the names have moved).
\end{boundary}

\section{Extended worked example: a two-object catalog, by hand}\label{ch08:sec:worked}

We now run Algorithm~\ref{ch08:alg:closure} by hand, round by round, on a
catalog small enough to trace completely yet rich enough to exercise every
branch: cross-object typing, non-injective collapse, frontier
deduplication in both its forms (accepted and queued), and termination.

\begin{example}[The catalog]\label{ch08:ex:workedcatalog}
Objects, in declared order:
$X=(\texttt{Left},(x_0,x_1))$ and $Y=(\texttt{Right},(y_0,y_1))$.
Generators, in declared order:
\begin{itemize}
\item $u$ (rank $0$), arrow $u:X\to Y$ with $u(x_0)=y_1$, $u(x_1)=y_0$
  (a ``crossing'' bijection into the other object);
\item $v$ (rank $1$), arrow $v:Y\to X$ with $v(y_0)=x_0$, $v(y_1)=x_0$
  (a collapse: everything to $x_0$).
\end{itemize}
Object ranks: $X\mapsto 0$, $Y\mapsto 1$. To keep the trace compact, name
the arrows that will appear:
$c_X:X\to X$ constant at $x_0$;
$c_Y:Y\to Y$ constant at $y_1$;
$c_{XY}:X\to Y$ constant at $y_1$;
$c_{YX}:Y\to X$ constant at $x_0$.
(The reader can verify as we go that these are exactly the composites that
arise.)
\end{example}

\emph{Seeding.} Push the empty word at $X$ (key $(0,(),0)$, denotation
$\id_X$) and the empty word at $Y$ (key $(0,(),1)$, denotation $\id_Y$).
Frontier: two entries; \emph{queued} $=\{\id_X,\id_Y\}$; \emph{reps}
empty.

\emph{Round 1.} Pop $()$ at $X$ --- least key. New denotation: accept
$\id_X$ with representative $()$. Expand: generators with source $X$: only
$u$. Extension $(u)$, denotation $u$ itself, key $(1,(0),0)$; not accepted,
not queued: push. Frontier: $\{()\,@\,Y,\ (u)\}$.

\emph{Round 2.} Pop $()$ at $Y$ (key $(0,(),1)$ beats $(1,(0),0)$ on
length). Accept $\id_Y$ as $()$. Expand: source-$Y$ generators: only $v$.
Push $(v)$, denotation $v$, key $(1,(1),1)$. Frontier:
$\{(u),\ (v)\}$.

\emph{Round 3.} Pop $(u)$ (rank sequence $(0)<(1)$). Accept $u$. Its target
is $Y$, so expand by $v$: word $(u,v)$, denotation
$v\circ u$. Compute: $x_0\mapsto u\mapsto y_1\mapsto v\mapsto x_0$;
$x_1\mapsto y_0\mapsto x_0$. So $v\circ u=c_X$. Push $(u,v)$, key
$(2,(0,1),0)$. Frontier: $\{(v),\ (u,v)\}$.

\emph{Round 4.} Pop $(v)$. Accept $v$. Target $X$; expand by $u$: word
$(v,u)$, denotation $u\circ v$: $y_0\mapsto x_0\mapsto y_1$;
$y_1\mapsto x_0\mapsto y_1$. So $u\circ v=c_Y$. Push $(v,u)$, key
$(2,(1,0),1)$. Frontier: $\{(u,v),\ (v,u)\}$.

\begin{center}
\begin{tabular}{@{}llll@{}}
\toprule
\multicolumn{4}{@{}l}{State after Round 4} \\
\midrule
Accepted denotation & Representative & Key & \\
\midrule
$\id_X$ & $()$ at $X$ & $(0,(),0)$ & \\
$\id_Y$ & $()$ at $Y$ & $(0,(),1)$ & \\
$u$     & $(u)$       & $(1,(0),0)$ & \\
$v$     & $(v)$       & $(1,(1),1)$ & \\
\midrule
Frontier & $(u,v)\mapsto c_X$ & $(2,(0,1),0)$ & pending \\
         & $(v,u)\mapsto c_Y$ & $(2,(1,0),1)$ & pending \\
\bottomrule
\end{tabular}
\end{center}

\emph{Round 5.} Pop $(u,v)$ --- rank sequence $(0,1)<(1,0)$. Accept $c_X$
with representative $(u,v)$: the first endomap in the table, and already a
non-obvious fact made canonical --- the least word producing the collapse
of $X$ routes through $Y$ and back. Expand: target $X$, generator $u$: word
$(u,v,u)$, denotation $u\circ c_X$: both points
$\mapsto x_0\mapsto y_1$, i.e.\ $c_{XY}$. Push, key $(3,(0,1,0),0)$.

\emph{Round 6.} Pop $(v,u)$. Accept $c_Y$ as $(v,u)$. Expand: target $Y$,
generator $v$: word $(v,u,v)$, denotation $v\circ c_Y$: both
$\mapsto y_1\mapsto x_0$, i.e.\ $c_{YX}$. Push, key $(3,(1,0,1),1)$.

\emph{Round 7.} Pop $(u,v,u)$. Accept $c_{XY}$ as $(u,v,u)$. Expand:
target $Y$, generator $v$: candidate $(u,v,u,v)$, denotation
$v\circ c_{XY}$: both points of $X\mapsto y_1\mapsto x_0$ --- that is
$c_X$, whose key is \emph{already accepted} (Round 5). Skip: no push. This
is the accepted-case prune of Lemma~\ref{ch08:lem:skip}, and the lemma's
guarantee is visible concretely: the retained representative $(u,v)$ is
shortlex-far below the skipped $(u,v,u,v)$.

\emph{Round 8.} Pop $(v,u,v)$. Accept $c_{YX}$ as $(v,u,v)$. Expand:
target $X$, generator $u$: candidate $(v,u,v,u)$, denotation
$u\circ c_{YX}=c_Y$ --- accepted in Round 6. Skip. Frontier now empty:
\textbf{terminate normally}.

\begin{center}
\begin{tabular}{@{}lll@{}}
\toprule
\multicolumn{3}{@{}l}{Final table: $8$ denotations (\texttt{denotation\_count} $=8$)} \\
\midrule
Denotation & Type & Representative \\
\midrule
$\id_X$   & $X\to X$ & $()$ at $X$ \\
$\id_Y$   & $Y\to Y$ & $()$ at $Y$ \\
$u$       & $X\to Y$ & $(u)$ \\
$v$       & $Y\to X$ & $(v)$ \\
$c_X$     & $X\to X$ & $(u,v)$ \\
$c_Y$     & $Y\to Y$ & $(v,u)$ \\
$c_{XY}$  & $X\to Y$ & $(u,v,u)$ \\
$c_{YX}$  & $Y\to X$ & $(v,u,v)$ \\
\bottomrule
\end{tabular}
\end{center}

Reconciliation with the theory: the closure size $8$ sits comfortably below
the Proposition~\ref{ch08:prop:finite} bound
$\sum_{i,j}2^2=4\cdot 4=16$; every representative is visibly the shortest
word denoting its arrow (and among equal-length candidates, none competed
here --- the catalog is too rigid); the peak frontier size was $2$,
illustrating Remark~\ref{ch08:rem:frontier}'s bound $F\le N$; and the two
skips in Rounds 7--8 are exactly the closure of the ``ping-pong'' orbit
$c_X\to c_{XY}\to c_X$ and $c_Y\to c_{YX}\to c_Y$: from length $3$ onward,
every longer word folds into one of the eight accepted denotations, which
is \emph{why} the algorithm halts. A reader wishing to internalize the
machine is invited to rerun the trace under the reversed generator order
$(v,u)$ and confirm: the same eight denotations, with the length-$1$ and
length-$3$ representatives unchanged as words but re-ranked, and the pop
order of Rounds 3--8 reshuffled accordingly
(Proposition~\ref{ch08:prop:orderrel}).

\begin{exercise}\label{ch08:exercise:third}
Extend the catalog of Example~\ref{ch08:ex:workedcatalog} with a third
generator $w$ (rank $2$) carrying the arrow $w:Y\to Y$ that swaps $y_0$ and
$y_1$. Before running the closure, predict: which new denotations become
reachable, which existing representatives change (careful: none of the
existing eight should --- adding a higher-ranked generator can only add
words that lose ties), and what \texttt{denotation\_count} becomes. Then
run the closure by hand and check. (The exercise rewards noticing that
$w\circ u$ ``uncrosses'' $u$, making the inclusion-like bijection
$X\to Y$ reachable at length $2$.)
\end{exercise}

\subsection{Lookups against the completed table}\label{ch08:sec:lookups}

The completed catalog of Example~\ref{ch08:ex:workedcatalog} also serves to
exercise the query surface and its error taxonomy, which is as much a part
of the construction's semantics as the table itself.

\begin{checkclass}[Error taxonomy of the query surface]\label{ch08:check:errors}
Three failure kinds are deliberately kept distinct, because they mean three
different things:
\begin{enumerate}
\item \emph{Malformed or foreign input} --- \texttt{ValidationError}. Raised
  by \texttt{evaluate} (hence by \texttt{normalize}) for an unknown
  generator label or an undeclared start object, and by the constructors
  for unordered or blank inputs. Meaning: the question itself is not a
  question about this catalog.
\item \emph{Ill-typed word} --- the composition type error of
  Chapter~\ref{ch04:main}, propagated unchanged through
  \texttt{evaluate}'s fold when a label's declared source fails to match
  the object reached so far. Meaning: the word is syntactically over this
  catalog's alphabet but does not type-check; it \emph{has no denotation},
  so no normal form is defined for it, and none is invented.
\item \emph{Unreachable arrow} --- \texttt{UnreachableDenotationError} from
  \texttt{representative\_for}. Meaning: the arrow is a perfectly valid
  typed arrow between declared objects, but the declared generators do not
  generate it. On a \emph{completed} table this verdict is sound and final
  (Theorem~\ref{ch08:thm:soundness}, completeness), which is exactly why
  the fail-closed rule matters: against a partial table the same lookup
  failure would be uninterpretable.
\end{enumerate}
\end{checkclass}

The third kind has a crisp instance in the worked catalog: the swap
endomap $\sigma:X\to X$, $x_0\leftrightarrow x_1$, is a legitimate arrow on
a declared object, yet the final table's only $X\to X$ entries are $\id_X$
and $c_X$. So \texttt{representative\_for($\sigma$)} raises
\texttt{UnreachableDenotationError} --- and that is a \emph{theorem} about
the catalog, not a shrug: since $v$'s collapse is the only way back into
$X$, every $X\to X$ word of positive length factors through $c_X$ and is
constant; no word denotes $\sigma$. The completed closure turns
``not found'' into ``provably not generated,'' which is the entire
difference between a cache miss and a mathematical statement.

\begin{exercise}\label{ch08:exercise:chain}
(i) Show that for a catalog whose generators are all endomaps of a single
declared object, Definition~\ref{ch08:def:shortlex} collapses to the
classical shortlex order on the free monoid over the ranked alphabet, the
third key component never differing. (ii) Construct, for each $n\ge 1$, a
single-object catalog with one generator whose closure has exactly $n$
denotations and whose longest representative has length $n-1$ (hint: a
cycle of length $n$; compare the worst-case length bound in
Remark~\ref{ch08:rem:cost}). (iii) Verify that in your construction the
frontier never holds more than one entry, so the peak $F=1$: the cost
bounds of Remark~\ref{ch08:rem:cost} are loose precisely when the
denotation graph is thin.
\end{exercise}

\section{Two permanent boundaries}\label{ch08:sec:boundaries}

\begin{boundary}[Not a rewrite-system claim]\label{ch08:bd:norewrite}
$\NF_G$ is \emph{denotational minimization over a completed finite
closure}. It is not --- and must never be quoted as --- a termination or
confluence result for a rewrite system, because \emph{no rewrite system is
defined here}. The distinction matters because the surrounding literature
creates the expectation: canonical forms for words modulo relations are
classically obtained by completing a set of rewrite rules and proving
termination and local confluence, in the tradition of string-rewriting
presentations of monoids and semigroups \cite{howie,pin}. The kernel
deliberately takes the other road: it never manipulates words by rules at
all; it \emph{evaluates} every word into the finite semantic domain of
typed total functions and minimizes over each semantic equivalence class
directly. Consequently: no claim is made that any set of TKS-inherited
equations, oriented in any way, terminates; no claim is made that any such
system is confluent; no critical-pair analysis exists; and no such analysis
is owed, because no result in this chapter depends on one. Should a future
lane define a rewrite presentation for inherited equations, its
termination and confluence obligations would be new ledger items, starting
from zero evidence --- nothing proved here transfers.
\end{boundary}

\begin{boundary}[Catalog identity, restated as a permanent rule]\label{ch08:bd:identity}
Restating Boundary~\ref{ch08:bd:catalogidentity} in its permanent,
cross-artifact form: normal forms compose with catalog identity and with
nothing else. A representative is a name \emph{within} a catalog, not a
name \emph{of} an arrow in the abstract. Every artifact that records a
normal form must record the catalog (or its digest) beside it; every
comparison of normal forms must first establish catalog equality; and
every catalog modification --- reordering, relabeling, adding or removing a
generator or object --- creates a \emph{new} catalog whose representatives
are incomparable with the old, even where
Proposition~\ref{ch08:prop:orderrel} guarantees the denotation sets agree.
The executable suite enforces the constructor-level half of this rule
(validated, ordered, hash-bindable inputs); the workflow-level half is a
discipline on consumers, recorded here as a permanent boundary so that no
future integration can cite this chapter as license for cross-catalog
comparison.
\end{boundary}

\section{Status summary}\label{ch08:sec:status}

\begin{center}
\begin{tabular}{@{}lp{5.0cm}p{4.6cm}@{}}
\toprule
Item & Where & Status \\
\midrule
Shortlex is a total well-order; keys injective &
Lemmas~\ref{ch08:lem:keyinjective}--\ref{ch08:lem:wellfounded} &
informal proof, this chapter \\
$\NF_G$ well defined &
Prop.~\ref{ch08:prop:nfwelldef} &
informal proof, this chapter \\
Dequeue monotonicity; skip domination; prefix domination &
Lemmas~\ref{ch08:lem:dequeue}, \ref{ch08:lem:skip}, \ref{ch08:lem:prefix} &
informal correctness argument (PO-TKS8-CORE-06) \\
Closure completeness and minimality &
Thm.~\ref{ch08:thm:soundness} &
informal correctness argument; formalization open; not in Lean verified
scope \\
Finiteness; termination; cost bounds &
Prop.~\ref{ch08:prop:finite}, Rem.~\ref{ch08:rem:cost} &
informal proof; bounds only \\
Fail-closed resource semantics &
Constr.~\ref{ch08:con:resources}, Rule~\ref{ch08:rule:failclosed} &
executable validation (resource regressions) \\
Order-relativity of representatives &
Prop.~\ref{ch08:prop:orderrel}, Ex.~\ref{ch08:ex:orderrel} &
informal proof; worked example \\
No rewrite-system claim; catalog identity required &
Boundaries~\ref{ch08:bd:norewrite}, \ref{ch08:bd:identity} &
permanent boundaries \\
\bottomrule
\end{tabular}
\end{center}

The chapter closes the mathematical development of the kernel's five
capabilities. What remains for the monograph's final part is the
engineering: how the implementation realizes these definitions
(Chapter~\ref{ch09:implementation}), how the bounded exhaustive suites and
regression fixtures exercise them (Chapter~\ref{ch10:testing}), what the
pinned Lean seed does and does not verify (Chapter 11), and where the
kernel's deliberate boundaries leave the wider TKS/CAS program
(Chapter 12). Nothing in this chapter licenses semantic transfer onto
canonical constructs; obligation PO-TKS8-CORE-08 remains \textsc{open / not
licensed}, and the shortlex machinery, like the lift before it, is offered
as substrate --- exact, replayable, fail-closed --- awaiting separately
gated review.
